\documentclass[12pt,a4paper]{article}

\usepackage[utf8]{inputenc}
\usepackage[margin=1in]{geometry}
\usepackage{amsmath, amssymb, amsthm}
\usepackage{graphicx}
\usepackage{xcolor}
\usepackage{caption}
\captionsetup[figure]{font=small}
\usepackage{float}
\usepackage{subcaption}
\usepackage{booktabs}
\usepackage{tabularx}
\usepackage{siunitx}
\usepackage{hyperref}
\usepackage{bm}
%\usepackage{natbib}
\usepackage{multirow}
\usepackage{makecell}

% ---------- PATH MACROS FOR FIGURES ----------
% convex combination plots
\newcommand{\convexBase}{../../../Results/Convex_combination}

\title{Simulated Time series Clustering in the $H \times C$ Plane Using Ordinal Pattern Features}
\author{Rasika Dilhani}
\date{\today}

\begin{document}
	
	\maketitle
	
\section{Purpose}
\noindent
The main purpose of this research is to use the simulated results from the previous study to support time series clustering. The simulation results show that the $H \times C$ representation captures two broad regimes of behavior: low‑entropy, high‑complexity dynamics and high‑entropy, low‑complexity dynamics. The diversity exhibited across these regimes confirms that ordinal pattern features provide meaningful discrimination between qualitatively different time series models. This separation in the $H \times C$ plane, together with the additional features derived from R\'enyi entropy, Tsallis entropy, and Fisher information, provides an enriched feature space that yields a more detailed description of each time series.

\medskip

\noindent
Based on these features, we apply unsupervised time series clustering. The goal is to determine whether time series generated by different mechanisms can be grouped correctly using only their ordinal pattern features and $H \times C$ plane representations. The clustering results are evaluated using visual inspection of the $H \times C$ plane and silhouette coefficients from the $k$-means algorithm.

	
\section{Methodology}{\label{sec:methodolpgy}}


In this experiment, we consider two deterministic signals (sine,logistic map), and a stochastic model  (ARMA, including two particular cases, namely AR and MA). Based on the outcomes of the simulation study, we fixed the model parameters. The following section describes the simulation setup and feature construction procedures used in our clustering experiment.

\textbf{Simulation and Feature Construction:}
\begin{itemize}
	\item \textbf{Embedding:} A fixed embedding dimension ($D = 4$) and time delay ($\tau = 1$) are used.
	
	\item \textbf{Model structure:} The models considered include $\mathrm{AR}(2)$, $\mathrm{MA}(2)$, $\mathrm{ARMA}(2,2)$, the sine function $f(t)$, and the logistic map $L_t$.
	
	\item \textbf{Simulation:} For each model and parameter set, $R = 50$ time series are simulated with lengths $n = 1000$ and $n = 10000$.
	
	\item \textbf{Features:} For each realization, the following quantities are computed:
	\begin{itemize}
		\item \emph{Entropy:} Shannon, Rényi, and Tsallis entropies, along with the Fisher information measure.
		\item \emph{Complexity:} Statistical complexity measures corresponding to each entropy.
		\item \emph{Variance:} Variances for each model to assess the stability of the extracted features.
		\item \emph{Confidence intervals:} Confidence intervals for entropy and complexity to further evaluate feature stability.
	\end{itemize}
	The experimental results is based on the following fourteen ordinal pattern features:
	\begin{itemize}
		\item $H_{\text{Shannon}}$, $H_{\text{Rényi}}$, $H_{\text{Tsallis}}$, $H_{\text{Fisher}}$,
		\item $C_{\text{Shannon}}$, $C_{\text{Rényi}}$, $C_{\text{Tsallis}}$, $C_{\text{Fisher}}$,
		\item Disequilibrium,
		\item $\mathrm{Var}(H_{\text{Shannon}})$, $\mathrm{Var}(H_{\text{Rényi}})$, $\mathrm{Var}(H_{\text{Tsallis}})$,
		% $\mathrm{Var}(H_{\text{Fisher}})$,
		\item $\mathrm{Var}(C_{\text{Shannon}})$.
	\end{itemize}
	
\subsection*{Information Theoretic Quantifiers}

For each ordinal pattern probability vector $P=\{p_\ell\}_{\ell=1}^{k}$, we evaluate a collection of 
global and local information measures.  
Given an estimator $\hat{P}=\{\hat{p}_\ell\}$ obtained from the empirical ordinal pattern frequencies, our goal is to study the behaviour of the entropy functionals $H(\hat{P})$ and their associated complexity measures.

\paragraph{Shannon entropy.}
The Shannon entropy is defined as
\[
H_{\mathrm{S}}(\hat{P})
= -\sum_{\ell=1}^{k} \hat{p}_\ell \log \hat{p}_\ell,
\]
and we use its normalized version $H_{\mathrm{S}}/H_{\max}$, where $H_{\max}$ corresponds to the 
uniform distribution.

\paragraph{Tsallis entropy.}
For $q \in \mathbb{R}\setminus\{1\}$, the Tsallis entropy is given by
\[
H^{\mathrm{T}}_{q}(\hat{P})
= \frac{1}{q-1}\left(1 - \sum_{\ell=1}^{k} \hat{p}_\ell^{\,q}\right),
\]
which reduces to Shannon entropy in the limit $q \to 1$.  
A normalized version is obtained by dividing by its value under the uniform distribution.

\paragraph{Rényi entropy.}
For $q>0$, $q\neq 1$, the Rényi entropy is defined as
\[
H^{\mathrm{R}}_{q}(\hat{P})
= \frac{1}{1-q}\log\!\left(\sum_{\ell=1}^{k} \hat{p}_\ell^{\,q}\right),
\]
and is likewise normalized using its maximum value at the uniform distribution.

\paragraph{Fisher Information Measure.}
To quantify local sensitivity in the ordinal pattern distribution and this is discussed by Felipe et.al.~\cite{Olivares2012b}, we adopt the discrete Fisher Information Measure proposed by Ferri et.al.~\cite{Ferri2009}:
\[
F(\hat{P})
= \frac{1}{2}\sum_{\ell=1}^{k-1}
\frac{(\hat{p}_{\ell+1}-\hat{p}_{\ell})^{2}}
{\hat{p}_{\ell+1}+\hat{p}_{\ell}}.
\]

This expression emphasizes differences between adjacent probabilities and therefore captures 
fine scale structure that entropy based measures may overlook.

\paragraph{Complexity.}
For each entropy type, the corresponding statistical complexity is computed as the product of the 
normalized entropy and the Jensen--Shannon divergence between $\hat{P}$ and the uniform 
distribution.  

	
	\item \textbf{Clustering:} Silhouette coefficients are used to assess the quality of clustering obtained via the $k$-means algorithm.
\end{itemize}

\noindent
For any stochastic model $\{X_t^{(S)}\}$ and deterministic model $\{X_t^{(D)}\}$, 
we define a convex combination as
\[
X_t(w) \;=\; w\,X_t^{(S)} + (1-w)\,X_t^{(D)}, 
\qquad 0 \le w \le 1,
\]
where the weight parameter $w$ controls the relative contribution of the stochastic and deterministic components.

\noindent
In our experiments, $X_t^{(S)} \in \{\mathrm{ARMA}(2,2), \mathrm{AR}(2), \mathrm{MA}(2)\}$ 
and $X_t^{(D)} \in \{\text{Logistic map } (L_t), \text{sine } f(t)\}$.
For each pair $(X_t^{(S)}, X_t^{(D)})$, the mapping
\[
w \longmapsto \bigl(H(X(w)),\, C(X(w))\bigr)
\]
traces the evolution of entropy and complexity coordinates as $w$ varies.

\noindent
The parameter settings for these models are summarized in Table~\ref{tab:parameters}. 
The experimental setup includes the following components:

\begin{enumerate}
	\item Three $\mathrm{ARMA}$ type models:
	\begin{itemize}
		\item $X_{\text{ARMA}} = \mathrm{ARMA}(2,2)$
		\item $X_{\text{AR}}   = \mathrm{AR}(2)$
		\item $X_{\text{MA}}   = \mathrm{MA}(2)$
	\end{itemize}
	
	\begin{table}[H]
		\centering
		\caption{Parameter settings for the $\mathrm{ARMA}$ type models used in the simulation study.}
		\label{tab:parameters}
		\begin{tabular}{lcc}
			\toprule
			\textbf{Model} & \textbf{AR coefficients} & \textbf{MA coefficients} \\ 
			\hline
			$\mathrm{ARMA}(2,2)$ & $(-0.8,\, 0.1)$ & $(-0.8,\, 0.1)$ \\ 
			\hline 
			$\mathrm{AR}(2)$     & $(0.1,\, 0.8)$ & -- \\ 
			\hline 
			$\mathrm{MA}(2)$     & -- & $(0.1,\, -0.8)$ \\ 
			\bottomrule
		\end{tabular}
	\end{table}
	
	\item A logistic map model defined as
	\[
	L_t = r L_{t-1}(1 - L_{t-1}),
	\]
	where $r = 3.8$ and the initial value $L_0 = 0.5$.
	
	\item A sine function defined as
	\[
	f(t) = \sin(2\pi \nu t),
	\]
	where $\nu = 0.05$ denotes the frequency.
	
	\item Convex combinations are constructed according to the general formula above. 
	These combinations were implemented in the simulation using weighted sums of the corresponding model outputs in \texttt{R}. 
	Table~\ref{tab:series_list} shows the representative model structure used to generate the series is illustrated as:
	
%	\begin{verbatim}
%		series_list <- list(
%		# Pure models
%		"ARMA(2,2)" = x_arma,
%		"AR(2)"     = x_ar2,
%		"MA(2)"     = x_ma2,
%		"Logistic"  = x_log,
%		"Sine"      = x_sin,
		
%		# Convex combinations of models
%		"ARMA+Sine(w=0.1)" = 0.1*x_arma + 0.9*x_sin,
%		"ARMA+Sine(w=0.2)" = 0.2*x_arma + 0.8*x_sin,
%		"ARMA+Sine(w=0.3)" = 0.3*x_arma + 0.7*x_sin,
		
%		# AR2 + Logistic
%		"AR2+Logistic(w=0.1)" = 0.1*x_ar2 + 0.9*x_log,
		
%		# AR2 + Sine
%		"AR2+Sine(w=0.8)" = 0.8*x_ar2 + 0.2*x_sin,
		
%		# MA2 + Sine
%		"MA2+Sine(w=0.4)" = 0.4*x_ma2 + 0.6*x_sin,
%		"MA2+Sine(w=0.6)" = 0.6*x_ma2 + 0.4*x_sin,
%		"MA2+Sine(w=0.8)" = 0.8*x_ma2 + 0.2*x_sin
%		)
%	\end{verbatim}
	
\begin{table}[H]
	\centering
	\caption{Pure models and convex combination time series used in the simulation.}
	\label{tab:series_list}
	\begin{tabular}{ll}
		\hline
		Name & Definition \\
		\hline
		ARMA(2,2) & $X_{\text{ARMA}}$ \\
		AR(2)     & $X_{\text{AR2}}$ \\
		MA(2)     & $X_{\text{MA2}}$ \\
		Logistic  & $X_{L_t}$ \\
		Sine      & $X_{f(t)}$ \\
		\hline
		$\mathrm{ARMA}$+Sine($w=0.1$) & $0.1\,X_{\text{ARMA}} + 0.9\,X_{f(t)}$ \\
		$\mathrm{ARMA}$+Sine($w=0.2$) & $0.2\,X_{\text{ARMA}} + 0.8\,X_{f(t)}$ \\
		$\mathrm{ARMA}$+Sine($w=0.3$) & $0.3\,X_{\text{ARMA}} + 0.7\,X_{f(t)}$ \\
		\hline
		$\mathrm{AR(2)}$+Logistic($w=0.1$) & $0.1\,X_{\text{AR2}} + 0.9\,X_{L_t}$ \\
		$\mathrm{AR(2)}$+Sine($w=0.8$) & $0.8\,X_{\text{AR2}} + 0.2\,X_{f(t)}$ \\
		\hline
		$\mathrm{MA(2)}$+Logistic($w=0.2$) & $0.2\,X_{\text{MA2}} + 0.8\,X_{L_t}$ \\
		$\mathrm{MA(2)}$+Logistic($w=0.7$) & $0.7\,X_{\text{MA2}} + 0.3\,X_{L_t}$ \\
		$\mathrm{MA(2)}$+Sine($w=0.4$) & $0.4\,X_{\text{MA2}} + 0.6\,X_{f(t)}$ \\
		$\mathrm{MA(2)}$+Sine($w=0.6$) & $0.6\,X_{\text{MA2}} + 0.4\,X_{f(t)}$ \\
		$\mathrm{MA(2)}$+Sine($w=0.8$) & $0.8\,X_{\text{MA2}} + 0.2\,X_{f(t)}$ \\
		\hline
	\end{tabular}
\end{table}
	
	
This formulation directly follows from the convex combination model $X_t(w)$, where each combination and weight value $w \in \{0.1, 0.2, \ldots, 0.8\}$ controls the mixture proportion between the chosen deterministic and stochastic components.
\end{enumerate}

\noindent
For each model and each weight value $w$, time series of length $n = 1000$ and $n=10000$ are simulated and the corresponding entropy and complexity measures are computed. 
An embedding dimension of $D = 4$ is applied throughout the experiments. 
All simulations are generated using the \texttt{Mersenne\_Twister} pseudorandom number generator.

Figure~\ref{fig:histogram_D4} shows 15 histograms generated from the above mentioned series.

\begin{figure}[H]
	\centering
	\includegraphics[width=\textwidth]{\convexBase/OP_Faceted_D4_n10000}
	\caption{Pattern histogram of selected time series for embedding dimension $D=4,$ time delay $\tau =1,$ and sequence length $10000$.}
	\label{fig:histogram_D4}
\end{figure}


\begin{figure}[H]
	\centering
	\includegraphics[width=\textwidth]{\convexBase/HC_TimeSeries_D4_n10000_Rep1}
	\caption{Fifteen systems and their points in the $H \times C$ plane with confidence intervals for embedding dimension $D=4,$ time delay $\tau =1,$ and sequence length $10000$.}
	\label{fig:timeseries_D4}
\end{figure}

Figure~\ref{fig:timeseries_D4} shows the lower right corner of the \(H \times C\) plane for embedding dimension \(D = 4\), with the time series highlighting the locations of the points \(\mathrm{MA}(2) + \sin(w = 0.8)\) and \(\mathrm{MA}(2)\).
The convex combination model uses weight \(0.8\) on the \(\mathrm{MA}(2)\) component and \(0.2\) on the sine function.
Although the sine function is deterministic, its behaviour is characteristic of stochastic or noise-dominated dynamics.
The ordinal patterns occur with nearly uniform probability, resulting in low structural organization despite high unpredictability.

We examined the correlations among the thirteen ordinal pattern features used in this study. 
The correlation matrix (Figure~\ref{fig:correlationfull_D4}) $(n = 10000)$ reveals several strong dependencies among the 13 ordinal pattern features. 
A detailed list and description of all thirteen features is given in Section~\ref{sec:methodolpgy}. 

The three entropy measures ($H_{\text{Shannon}}$, $H_{\text{Rényi}}$, $H_{\text{Tsallis}}$) exhibit extremely strong positive correlations ($r = 0.99$ for all pairs), indicating that they capture nearly identical distributional information. 
In contrast, $H_{\text{Fisher}}$ shows moderate negative correlations with these entropies, with values ranging from $r = -0.50$ (vs.\ $H_{\text{Tsallis}}$) to $r = -0.69$ (vs.\  $H_{\text{Shannon}}$).  

A similar pattern appears among the statistical complexity measures ($C_{\text{Shannon}}$, $C_{\text{Rényi}}$, $C_{\text{Tsallis}}$), which show very strong positive correlations 
($r = 0.99$ for all pairs). 
Strong positive associations are also observed between the complexity measures and $C_{\text{Fisher}}$ ($r = 0.87$–$0.94$).  

Entropy measures and disequilibrium exhibit extremely strong negative correlations 
($r = -0.99$ for $H_{\text{Shannon}}$ vs.\ Disequilibrium, 
$r = -0.90$ for $H_{\text{Rényi}}$, 
$r = -0.95$ for $H_{\text{Tsallis}}$). 
Entropy increases when the pattern distribution becomes more random, whereas disequilibrium increases when the distribution becomes more uneven; because these behaviours oppose each other, the two measures are strongly negatively correlated.

%\begin{figure}[H]
%	\centering
%	\includegraphics[width=\textwidth]{\convexBase/Correlation}
%	\caption{Correlation matrix for sequence length $10000$.}
%	\label{fig:correlation_D4}
%\end{figure}

\begin{figure}[H]
	\centering
	\includegraphics[width=\textwidth]{\convexBase/Correlation_full}
	\caption{Correlation matrix for embedding dimension $D=4,$ time delay $\tau =1,$ and sequence length $10000$.}
	\label{fig:correlationfull_D4}
\end{figure}

Variance-based features also display high internal correlation, particularly between 
$\mathrm{Var}(H_{\text{Rényi}})$ and $\mathrm{Var}(H_{\text{Tsallis}})$ ($r = 0.93$), 
and between $\mathrm{Var}(H_{\text{Shannon}})$ and both of these variances ($r = 0.93$). 
In contrast, $\mathrm{Var}(C_{\text{Shannon}})$ shows only weak correlations with most other features ($|r| \le 0.35$). 
Fisher entropy, and Fisher complexity also exhibits weak associations with the entropy variances ($|r| \le 0.19$).  

Overall, the matrix indicates substantial multicollinearity among entropy and complexity based features, while Fisher-related and variance-based descriptors provide more independent information for clustering.  
Disequilibrium also shows only weak to moderate positive correlations with the entropy variances 
($r = 0.30$–$0.37$), reinforcing that it behaves largely independently from both entropy and variance-based quantities.

Also, we checked the variability of the features across different classes of ARMA models and convex combination models. 
%Figure~\ref{fig:boxplot1} presents the boxplots of the features for each class. The results are quite clear; however, the variance of Fisher measure is on a substantially different scale, making it difficult to assess the variability of the remaining features.

%\begin{figure}[H]
%	\centering
%	\includegraphics[width=\textwidth]{\convexBase/Boxplots_Normal}
%	\caption{Comparative Boxplots of the features for sequence length $10000$.}
%	\label{fig:boxplot1}
%\end{figure}

To better evaluate how the features vary by class, and noting that all measures are nonnegative (it was confirmed by the summary statistics), we constructed boxplots on a logarithmic scale. This transformation allows for improved visualization of feature variability, particularly for measures spanning a wide range of values (see Figure~\ref{fig:boxplot2}).

\begin{figure}[H]
	\centering
	\includegraphics[width=\textwidth]{\convexBase/Boxplots_LogScale}
	\caption{Comparative Boxplots of the features by logarithmic scale for sequence length $10000$.}
	\label{fig:boxplot2}
\end{figure}

It is clear that some features vary across several orders of magnitude, while others remain much more concentrated. 
In particular, measures such as the variance based features show wide vertical spreads within certain classes, whereas entropy and complexity measures tend to cluster more tightly, indicating greater stability across models.

\section{Time Series Clustering}
Clustering is an unsupervised learning process that involves organizing a collection of patterns into similar groups based on their similarity. 
The term \textit{unsupervised} refers to constructing a model or partitioning observations into groups directly from the data, without prior labels. 
The term \textit{learning} indicates that the statistician calibrates models to the available data. 
Once clusters or groups are defined, supervised learning can be applied, where \textit{supervised} means that the group membership is known a priori. 
Clustering is one of the most important problems in data mining and has numerous applications across various fields such as the natural sciences, the medical sciences, economics, marketing. In the context of time series analysis, clustering groups time series with similar temporal patterns, dynamics, or statistical properties.
Once clusters are obtained, it is generally useful to describe and interpret each group. 
Identifying compact and well-separated clusters in data, however, remains a challenging task.


As discussed in the book note in "Applied Multivariate Statistical Analysis"~\cite{HardleSimarFengler2024}, clustering reveals latent structure in high-dimensional data, making it fundamental for exploratory analysis.
In this research work, we focus on simulated time series clustering using the $k$-means algorithm and validate the results using silhouette coefficients.

\subsection{Simulation study to assess the clustering procedure}

To evaluate the performance of the clustering approach, time series were simulated from a collection of fifteen time series, including the convex combination models based on $\mathrm{ARMA}(2,2)$, $\mathrm{AR}(2)$, $\mathrm{MA}(2)$, Logistic, Sine, and their hybrid extensions as described in the Table~\ref{tab:series_list}. 
For each model class, multiple series of lengths 1000 and 10000 were generated, and the clustering procedure was applied. 
This entire process was repeated 50 times in order to evaluate the stability of the clustering outcomes.  

For each repetition, we recorded the number of correctly identified clusters, where a correct cluster is defined as one containing all series generated from the same underlying process. 

The clustering is based on the following fourteen ordinal‑pattern features:
\begin{itemize}
	\item $H_{\text{Shannon}}$, $H_{\text{Rényi}}$, $H_{\text{Tsallis}}$, $H_{\text{Fisher}}$,
	\item $C_{\text{Shannon}}$, $C_{\text{Rényi}}$, $C_{\text{Tsallis}}$, $C_{\text{Fisher}}$,
	\item Disequilibrium,
	\item $\mathrm{Var}(H_{\text{Shannon}})$, $\mathrm{Var}(H_{\text{Rényi}})$, $\mathrm{Var}(H_{\text{Tsallis}})$, 
	%$\mathrm{Var}(H_{\text{Fisher}})$,
	\item $\mathrm{Var}(C_{\text{Shannon}})$.
\end{itemize}

Both visual inspection of the clustering patterns in the $H \times C$ plane and the silhouette based evaluation are used to compare and validate the algorithmic results. Moreover, we use density based (DBSCAN) clustering method to further verify our results. 

\subsection{Silhouette-Based Cluster Evaluation}
The silhouette coefficient provides a quantitative measure of how well each time series is assigned to its cluster. 
For a given observation $i$, the silhouette value is defined as
\[
s(i) = \frac{b(i) - a(i)}{\max\{a(i),\, b(i)\}},
\]
where $a(i)$ denotes the average distance between $i$ and all other points in the same cluster, and $b(i)$ is the minimum average distance between $i$ and all points in any other cluster. 
The silhouette score therefore lies in the interval $[-1,1]$, with values close to $1$ indicating well-separated clusters, values near $0$ suggesting ambiguous assignments, and negative values indicating potential misclassification.

In this study, the silhouette index is used to assess the internal consistency of the clustering results obtained from the ordinal pattern feature space. 
For each repetition of the experiment, we compute the average silhouette width across all time series, providing a global measure of cluster quality. 
Higher silhouette values correspond to more coherent groupings of time series generated from the same underlying model, and thus serve as an objective complement to the visual inspection of the $H \times C$ plane.

\subsection{Density-Based Clustering via DBSCAN}

To further validate the clustering structure, we employ the Density-Based Spatial Clustering of Applications with Noise (DBSCAN) algorithm, a widely used density-based method originally introduced by \cite{Ester1996}. 
DBSCAN is based on a simple idea: clusters appear in parts of the feature space where many points lie close together, while regions with only a few points are treated as gaps or noise. 
This makes DBSCAN different from methods such as $k$-means, which assume that the number of clusters is known in advance.

DBSCAN uses two parameters: the neighbourhood radius $\varepsilon$ and the minimum number of points needed to form a dense region, denoted $\text{minPts}$. 
For any point $p$, the $\varepsilon$-neighbourhood is defined as
\[
\mathcal{N}_\varepsilon(p) = \{q : d(p,q) \le \varepsilon\}.
\]

A point is classified as:
\[
\text{core point if } |\mathcal{N}_\varepsilon(p)| \ge \text{minPts},
\]
border point if it is close to a core point but does not meet the density requirement, noise point otherwise

Clusters are formed by linking core points that are close to each other.  
If two core points can be reached through a chain of other core points, they belong to the same cluster.  
Border points are attached to the nearest core region, while noise points remain unassigned.

This density-based view allows DBSCAN to detect clusters of irregular shape and to separate meaningful structure from isolated observations. 
It is therefore well suited for the ordinal‑pattern feature space, where the geometry of the data may be complex and not well captured by distance‑based partitioning methods.

In our analysis, the value of $\varepsilon$ is chosen using the $k$‑nearest neighbour distance curve, selecting the point where the curve shows a clear ``elbow''. 
The parameter $\text{minPts}$ is set to the dimensionality of the feature space by adding one ($\text{minPts} = 14$). 
DBSCAN is applied to each repetition of the experiment to check whether the density-based grouping agrees with the clusters obtained from the primary algorithmic approach.


\section{Results and Discussion}
This section discusses the clustering results. Based on the poor performance of the clustering algorithm for the sample size 1000, we focus on the sample size 10000 when assessing the clustering behaviour. Figures~\ref{fig:Shannon scatter},\ref{fig:Renyi scatter}, \ref{fig:Tsallis scatter}, and~\ref{fig:Fisher scatter} show scatter plots in the $H \times C$ plane for each of the 15 time series models considered in our experimental setup. From these plots, it is visually evident that some classes can be distinguished in the \(H \times C\) plane. However, we formally confirm these visual findings using cluster analysis.

\begin{figure}[H]
	\centering
	\includegraphics[width=\textwidth]{\convexBase/HC_Shannon_D4_n1000}
	\caption{Shannon Entropy and Complexity Scatter plot for embedding dimension $D=4,$ time delay $\tau =1,$ and sequence length $1000$.}
	\label{fig:Shannon scatter1000}
\end{figure}

\begin{figure}[H]
	\centering
	\includegraphics[width=\textwidth]{\convexBase/HC_Shannon_D4_n10000}
	\caption{Shannon Entropy and Complexity Scatter plot for embedding dimension $D=4,$ time delay $\tau =1,$ and sequence length $10000$.}
	\label{fig:Shannon scatter}
\end{figure}

\begin{figure}[H]
	\centering
	\includegraphics[width=\textwidth]{\convexBase/HC_Renyi_D4_n10000}
	\caption{R\'eny Entropy and Complexity Scatter plot for embedding dimension $D=4,$ time delay $\tau =1,$ and sequence length $10000$.}
	\label{fig:Renyi scatter}
\end{figure}

\begin{figure}[H]
	\centering
	\includegraphics[width=\textwidth]{\convexBase/HC_Tsallis_D4_n10000}
	\caption{Tsallis Entropy and Complexity Scatter plot for embedding dimension $D=4,$ time delay $\tau =1,$ and sequence length $10000$.}
	\label{fig:Tsallis scatter}
\end{figure}

\begin{figure}[H]
	\centering
	\includegraphics[width=\textwidth]{\convexBase/HC_Fisher_D4_n10000}
	\caption{Fisher Information and Complexity Scatter plot for embedding dimension $D=4,$ time delay $\tau =1,$ and sequence length $10000$.}
	\label{fig:Fisher scatter}
\end{figure}

Visual inspection of the $H \times C$ plane indicates that the clusters become clearly separable once the sample size is sufficiently large. 
To confirm these visually observed patterns, we apply the $k$-means clustering algorithm and evaluate the resulting partition using silhouette coefficients.

Looking at the elbow plot (Figure~\ref{fig:elbow}) and silhouette score plot (Figure~\ref{fig:silhouette}), it is clear that we can distinguish 11 clusters. 

\begin{figure}[H]
	\centering
	\includegraphics[width=\textwidth]{\convexBase/Elbow}
	\caption{Elbow plot for embedding dimension $D=4,$ time delay $\tau =1,$ and sequence length $10000$.}
	\label{fig:elbow}
\end{figure}

\begin{figure}[H]
	\centering
	\includegraphics[width=\textwidth]{\convexBase/Silhouette}
	\caption{Silhouette coefficients plot for embedding dimension $D=4,$ time delay $\tau =1,$ and sequence length $10000$.}
	\label{fig:silhouette}
\end{figure}


The clustering results are summarized in Figures~\ref{fig:cluster_sizes} and \ref{fig:silhouette_coef}.  

\begin{figure}[H]
	\centering
	\includegraphics[width=\textwidth]{\convexBase/Cluster_Sizes}
	\caption{Cluster sizes for embedding dimension $D=4,$ time delay $\tau =1,$ and sequence length $10000$.}
	\label{fig:cluster_sizes}
\end{figure}

Figure~\ref{fig:cluster_sizes} illustrates the distribution of cluster sizes obtained from the cluster analysis with $k = 11$. 
A substantial variability in cluster sizes is observed, with some clusters containing 25 observations, while others include between 100 and 200 observations. 
This imbalance suggests that certain ordinal pattern occurs more frequently across the simulated models, whereas others correspond to more specialized or less common dynamical behaviours.

Based on the cluster assignments, the dominant model class within each cluster can be identified. 
Cluster~1 is primarily associated with the Sine process, while Clusters~2, 3, 5, 6, and 10 are dominated by $\mathrm{ARMA}(2,2)$, $\mathrm{MA2+Logistic}(0.2)$, $\mathrm{ARMA+Sine}(w=0.2)$, $\mathrm{ARMA+Sine}(w=0.3)$, and $\mathrm{MA2+Sine}(0.4)$, respectively. 
Cluster~4 groups together the $\mathrm{AR}(2)$ and $\mathrm{AR2+Sine}(0.8)$ processes, whereas Cluster~8 corresponds to the Logistic map and its hybrid extension with $\mathrm{AR}(2)$. 
Cluster~9 aggregates several moving‑average‑based hybrid processes, including $\mathrm{MA}(2)$, $\mathrm{MA2+Logistic}(0.7)$, $\mathrm{MA2+Sine}(0.6)$, and $\mathrm{MA2+Sine}(0.8)$.

Clusters~8 and 11 are not well separated, as they split the $\mathrm{ARMA+Sine}(0.1)$ series into two halves. 
This indicates that borderline hybrid cases such as $\mathrm{ARMA+Sine}$ naturally lie between neighbouring dynamical regimes, making them more difficult to assign cleanly to a single cluster.

These findings demonstrate that ordinal pattern features provide a meaningful and interpretable separation of model classes. 
Purely stochastic models, deterministic systems, and convex‑combination processes tend to form distinct groups in the feature space.
 
A summary of the dominant model associated with each cluster is provided in Table~\ref{tab:cluster_dominant_models}.

\begin{table}[H]
	\centering
	\caption{Cluster labels and dominant models}
	\label{tab:cluster_dominant_models}
	\begin{tabular}{cl}
		\hline
		Cluster & Dominant Models \\
		\hline
		1  & Sine \\
		2  & $\mathrm{ARMA(2,2)}$ \\
		3  & $\mathrm{MA(2)}$+Logistic($w=0.2$) \\
		4  & $\mathrm{AR(2)}$, $\mathrm{AR(2)}$+Sine($w=0.8$) \\
		5  & $\mathrm{ARMA}$+Sine($w=0.2$)  \\
		6  & $\mathrm{ARMA}$+Sine($w=0.3$)\\
		7  & Logistic, $\mathrm{AR(2)}$+Logistic ($w=0.1$) \\
		8  & $\mathrm{ARMA}$+Sine($w=0.1$) \\
		9  & $\mathrm{MA(2)}$, $\mathrm{MA(2)}$+Logistic($w=0.7$), $\mathrm{MA(2)}$+Sine($w=0.6$), $\mathrm{MA(2)}$+Sine($w=0.8$)  \\
		10 & $\mathrm{MA(2)}$+Sine($w=0.4$) \\
		11 & $\mathrm{ARMA}$+Sine($w=0.1$) \\
		\hline
	\end{tabular}
\end{table}

\begin{figure}[H]
	\centering
	\includegraphics[width=\textwidth]{\convexBase/Silhouette_Final_k}
	\caption{Silhouette coefficients for embedding dimension $D=4,$ time delay $\tau =1,$ and sequence length $10000$.}
	\label{fig:silhouette_coef}
\end{figure}

\begin{table}[H]
	\centering
	\caption{Cluster size and average silhouette width}
	\label{tab:cluster_silhouette}
	\begin{tabular}{rcc}
		\hline
		Cluster & Size & \texttt{ave.sil.width} \\
		\hline
		1  & 50  & 1.00 \\
		2  & 50 & 0.83 \\
		3  & 50  & 0.88 \\
		4  & 100 & 0.68 \\
		5  & 50  & 0.76 \\
		6  & 50  & 0.57 \\
		7  & 100  & 0.62 \\
		8  & 25  & 0.49 \\
		9  & 200  & 0.76 \\
		10 & 50  & 0.76 \\
		11 & 25  & 0.55 \\
		\hline
	\end{tabular}
\end{table}


The silhouette plot in Figure~\ref{fig:silhouette_coef} further supports the quality of the clustering. 
Most clusters exhibit positive silhouette widths, indicating that observations are closer to their own cluster center than to others, and the overall silhouette average for $k=12$ is the highest among all tested values. 
A few clusters have lower silhouette values, meaning some overlap exists $(\mathrm{ARMA(2,2)})$, but none of the clusters show a large number of negative silhouettes. This is also validate with the Table~\ref{tab:cluster_silhouette}  
This indicates that the clustering solution remains stable and most series are well assigned.
Taken together, the cluster sizes, the model to cluster mapping, and the silhouette coefficients provide strong evidence that the fourteen ordinal pattern features yield a coherent and informative partition of the simulated time series.

To further interpret the clustering results, a two dimensional PCA projection was used to visualize the distribution of the time series in the ordinal pattern features. 
This provides an intuitive representation of how the clusters (Figure~\ref{fig:pca_clusters}) identified by the $k$‑means algorithm are separated, and how these data points compare with the true underlying model classes(Figure~\ref{fig:pca_classes}).

\begin{figure}[H]
	\centering
	\includegraphics[width=\textwidth]{\convexBase/PCA_Clusters}
	\caption{PCA Clusters for embedding dimension $D=4,$ time delay $\tau =1,$ and sequence length $10000$.}
	\label{fig:pca_clusters}
\end{figure}

Figure~\ref{fig:pca_clusters} shows the PCA projection of all time series coloured by their assigned cluster labels ($k = 12$). 
The clusters form visibly distinct groups in the two–dimensional PCA space, though some clusters lie closer together, reflecting mild overlap in their ordinal‑pattern features. 
Overall, the plot demonstrates that the $k$‑means partition captures major structural differences in the feature space and provides a low dimensional view of cluster separation.

\begin{figure}[H]
	\centering
	\includegraphics[width=\textwidth]{\convexBase/PCA_Classes}
	\caption{PCA Classes for embedding dimension $D=4,$ time delay $\tau =1,$ and sequence length $10000$.}
	\label{fig:pca_classes}
\end{figure}

Figure~\ref{fig:pca_classes} presents the PCA projection coloured according to the true generating model of each time series. 
The true classes also appear well separated in the first two principal components, with nonlinear models such as the Logistic and Sine processes forming distinct regions, while convex combination ARMA–Sine and AR2–Logistic processes occupy characteristic intermediate positions. 
This visual structure supports the conclusion that the ordinal pattern features provide strong discriminatory power, since most model classes form identifiable clusters even before applying any clustering algorithm.

Further, the DBSCAN analysis confirms that it is possible to differentiate 9 clusters from the 15 simulated classes. First, the optimal value of $\varepsilon$ for the DBSCAN algorithm was determined using the elbow plot shown in Figure~\ref{fig:dbscan_elbow}. Based on the location of the elbow point, the value of $\varepsilon = 0.8$ was selected for the clustering analysis.

\begin{figure}[H]
	\centering
	\includegraphics[width=\textwidth]{\convexBase/kNN_D4}
	\caption{DBSCAN Elbow plot.}
	\label{fig:dbscan_elbow}
\end{figure}

After selecting the optimal parameter, the cluster structure was examined using the \texttt{fviz\_cluster} representation shown in Figure~\ref{fig:dbscan_fviz}. A similar clustering pattern can also be observed in the PCA projection presented in Figure~\ref{fig:dbscan_pca}. Both visualisations indicate a clear separation between several groups in the ordinal pattern feature space.

\begin{figure}[H]
	\centering
	\includegraphics[width=\textwidth]{\convexBase/fviz_cluster_D4}
	\caption{DBSCAN cluster structure using the fviz cluster.}
	\label{fig:dbscan_fviz}
\end{figure}

\begin{figure}[H]
	\centering
	\includegraphics[width=\textwidth]{\convexBase/DBSCAN_PCA_D4}
	\caption{DBSCAN PCA projection.}
	\label{fig:dbscan_pca}
\end{figure}

The resulting cluster assignments are summarised in Table~\ref{tab:cluster_grouped}. The results show that several related models are grouped into the same cluster, while others form distinct isolated groups. In particular, the models based on $\mathrm{MA}(2)$ dynamics and several of its hybrid variations are consistently assigned to Cluster~2, indicating strong structural similarity in the ordinal pattern feature space. In contrast, the Logistic and Sine models each form their own individual clusters, demonstrating that these nonlinear dynamics are clearly distinguishable from the remaining classes. 

The ARMA+Sine mixtures exhibit a progressive separation as the weight parameter increases, producing three different clusters for $w=0.1$, $w=0.2$, and $w=0.3$. This behaviour suggests that the ordinal pattern features are sensitive to gradual changes in the underlying signal composition. Overall, the DBSCAN results demonstrate that the entropy complexity representation preserves meaningful structural differences between the simulated time series models and allows most classes to be separated successfully.

\begin{table}[ht]
	\centering
	\caption{Classes grouped according to cluster assignments}
	\begin{tabular}{cl}
		\hline
		\textbf{Cluster} & \textbf{Class} \\
		\hline
		1 & ARMA(2,2), AR(2), AR2+Sine(w=0.8) \\
		2 & MA(2), MA2+Logistic(w=0.7), MA2+Sine(w=0.4), \\
		& MA2+Sine(w=0.6), MA2+Sine(w=0.8) \\
		3 & Logistic \\
		4 & Sine \\
		5 & ARMA+Sine(w=0.1) \\
		6 & ARMA+Sine(w=0.2) \\
		7 & ARMA+Sine(w=0.3) \\
		8 & AR2+Logistic(w=0.1) \\
		9 & MA2+Logistic(w=0.2) \\
		\hline
	\end{tabular}
	\label{tab:cluster_grouped}
\end{table} 

Overall, the ordinal pattern features provide a much clearer and more informative representation of the time series than the clustering algorithm alone. 
The structure visible in the $H \times C$ plane allows the fifteen true model classes to be distinguished almost entirely by visual inspection, demonstrating that the ordinal pattern feature space itself contains strong separability. 
The $k$-means algorithm, supported by silhouette coefficients, confirms this visual evidence by recovering a similar cluster structure. 
These results indicate that the entropy complexity representation offers substantial insight into the underlying dynamics, and that the ordinal pattern features are highly effective for identifying and differentiating the simulated time series models. 
The next step is to extend this analysis to real world time series, where the same ordinal pattern framework can be used to evaluate whether comparable clustering behaviour arises outside a controlled simulation setting.
	
	
	\bibliographystyle{acm}
	\bibliography{../../BearingFaultDiagnosis}
	
\end{document}